\section{本章小结}
\label{sec:lln-summary}

本章研究独立随机变量的部分和$S_n$，并区分两个问题：未归一化的累计$\sum_kX_k$是否收敛，以及平均$\bar{X}_n$是否稳定在期望附近。两者用的是同一组工具，各工具解决的环节不同。

\paragraph{部分和的两个基本估计}
独立性使中心化交叉项的期望为零，方差可加，$\operatorname{Var}(S_n)=\sum_k\sigma_k^2$；部分和的散布因此按$\sqrt{n}$增长，样本平均按$1/\sqrt{n}$收窄。Chebyshev不等式只控制固定终点。Kolmogorov最大不等式按首次越过阈值的时刻拆分事件，利用之后增量的独立性与零均值，用同一个终点方差控制整段路径的最大偏离。它需要独立与中心化，两者缺一不可。

\paragraph{随机级数}
最大不等式控制尾段的振荡，于是独立中心化、方差可和的级数几乎必然收敛，同时在$L^2$中收敛。有非零均值时还要单独检查均值级数。方差或均值被罕见大值破坏时，在固定水平截断：大项由超阈值概率可和与Borel–Cantelli引理排除，截断均值由数项级数控制，截断波动由方差可和控制。三系列定理说明这三个条件也是必要的；随机符号级数$\sum_ka_k\varepsilon_k$几乎必然收敛当且仅当$\sum_ka_k^2<\infty$。

\paragraph{两种大数律}
弱大数律回答每个固定时刻偏差的概率。有限方差时Chebyshev不等式给出有限样本的界$\sigma^2/(n\varepsilon^2)$；只有一阶矩时，先固定截断水平、令样本量增大，再放大截断水平。强大数律回答几乎每条路径的最终行为。Kronecker引理把加权级数$\sum_k(X_k-\mu_k)/k$的收敛逐样本地翻译成平均的收敛；对可积的独立同分布序列，按项数$k$截断，使大项只出现有限次、截断方差的加权和被一阶矩控制。可积对强律既充分又必要，Cauchy分布是反例。

\paragraph{使用时核对的条件}
频率、函数平均与样本矩都是独立同分布变量的平均。使用前确认被平均的事件或函数在观测之前固定、相应的期望绝对可积；样本方差的几乎必然收敛只需有限二阶矩。对依赖数据挑选的事件族，单个事件的结论不能直接推广。

大数律确定了平均的极限位置，却没有说明平均在极限附近怎样波动。把偏差$\bar{X}_n-\mu$放大$\sqrt{n}$倍以后，其分布趋于什么形状，需要先建立只比较分布的收敛概念，这是第~\ref{chap:weak-convergence-of-distributions}章与第~\ref{chap:central-limit-theorem}章的内容。
